% Druckvorlage zur Übung Schulaufgabe 1 (Rekursive Datenstrukturen, Informatik 12)
% Übersetzen: erst pdflatex abbildungen.tex, dann pdflatex vorlagen.tex (oder ./bauen.sh)
\documentclass[11pt,a4paper]{article}
\usepackage[utf8]{inputenc}
\usepackage[T1]{fontenc}
\usepackage{lmodern}
\usepackage[ngerman]{babel}
\usepackage[top=2cm,bottom=2cm,left=2cm,right=2cm]{geometry}
\usepackage{graphicx}
\usepackage{array}
\usepackage{tabularx}
\usepackage{fancyhdr}
\pagestyle{fancy}\fancyhf{}
\lhead{\textbf{Informatik 12} -- Übung zur 1. Schulaufgabe}
\rhead{Vorlagen zum Ausdrucken}
\cfoot{edu-mrh.de \quad Seite \thepage}
\setlength{\parindent}{0pt}
\newcommand{\abb}[1]{\par\noindent\includegraphics[width=\linewidth,page=#1]{abbildungen.pdf}\par}
\newcommand{\schreibzeile}{\rule[-2.8mm]{0pt}{8mm}}
\newcommand{\kaestchen}{\framebox[9mm]{\rule[-2.5mm]{0pt}{8mm}}}
\begin{document}

\section*{Aufgabe 1b: Pizzadienst}
\renewcommand{\arraystretch}{1}
\begin{tabularx}{\linewidth}{|p{3cm}|p{3cm}|X|}\hline
\textbf{Aufruf} & \textbf{Rückgabe} & \textbf{Inhalt danach (vorne $\rightarrow$ hinten)}\\\hline
\texttt{add(P1)}\schreibzeile & & \\\hline
\texttt{add(P2)}\schreibzeile & & \\\hline
\texttt{peek()}\schreibzeile & & \\\hline
\texttt{poll()}\schreibzeile & & \\\hline
\texttt{add(P3)}\schreibzeile & & \\\hline
\texttt{poll()}\schreibzeile & & \\\hline
\texttt{poll()}\schreibzeile & & \\\hline
\texttt{poll()}\schreibzeile & & \\\hline
\texttt{isEmpty()}\schreibzeile & & \\\hline
\end{tabularx}

\section*{Aufgabe 2a und 2b: Klassendiagramm der Playlist}
\abb{1}

\newpage
\section*{Aufgabe 2d: \texttt{liste.gesamtdauer()}}
\abb{2}
\section*{Aufgabe 2e: \texttt{liste.entfernen(B)}}
\abb{3}

\newpage
\section*{Aufgabe 3: \texttt{einfuegenAn} (Sequenzdiagramm als Hilfe)}
\abb{4}
\vspace{4mm}
{\ttfamily public ListenElement einfuegenAn(int position, DatenElement d) \{}\\[2mm]
\hspace*{8mm}\kaestchen\ \kaestchen\ \kaestchen\ \kaestchen\ \kaestchen\ \kaestchen\ \kaestchen\\[2mm]
{\ttfamily \}}

\section*{Aufgabe 5: Bildergalerie}
\abb{5}
\vspace{2mm}
\begin{tabularx}{\linewidth}{|X|X|X|X|}\hline
\textbf{Knoten} & \textbf{Attribut} & \textbf{alter Wert} & \textbf{neuer Wert}\\\hline
\schreibzeile & & & \\\hline
\schreibzeile & & & \\\hline
\schreibzeile & & & \\\hline
\schreibzeile & & & \\\hline
\schreibzeile & & & \\\hline
\end{tabularx}

\end{document}
